\uuid{2oHj}
\titre{Fonction croissante/décroissante}
\auteur{bodin ; hanani ; mzari}
\organisation{exo7}
\licence{CC BY-NC-SA 4.0}
\datecreate{2020-01-20}
\niveau{L1}
\module{Analyse}
\chapitre{Continuité, limite et étude de fonctions réelles}
\sousChapitre{Autre}
\difficulte{2}
\theme{Notion de fonctions, fonction croissante, décroissante}
\exo7qcmid{lille-1.continuite.004}

\contenu{
\begin{question}
\qclassification{123.00}
\qtags{motcle=fonction croissante/décroissante}
Quels arguments sont valables pour montrer que $f : \Rr \to \Rr$ est décroissante ?
\begin{answers}
    \bad{On a $x \le y$ qui implique $f(x) \le f(y)$.}

    \good{On a $x \le y$ qui implique $f(x) \ge f(y)$.}

    \good{On a $x \ge y$ qui implique $f(x) \le f(y)$.}

    \bad{On a $x \ge y$ qui implique $f(x) \ge f(y)$.}
\end{answers}
\begin{explanations}
Une fonction $f$ est décroissante si $x \le y$ implique $f(x) \ge f(y)$.
Ce qui peut aussi s'écrire $x \ge y$ qui implique $f(x) \le f(y)$.
Autrement dit $f$ inverse le sens des inégalités.
\end{explanations}
\end{question}
}
